

Il est
évidemment nécessaire pour cela que $X$ soit caténaire, et suffisant que $X$ soit biéquidimensionne (EGA $0_{IV}$ 14.3.3). Cependant il existe des schémas noethériens de dimension finie, caténaires et équidimensionnels, n'ayant pas de telles fonctions $D$. Pour le voir on peut modifier l'exemple (EGA IV 5.6.11) de la façon suivante : prenons deux exemplaires $Y_1$ et $Y_2$ du spectre de l'anneau $E$ de loc. cit. et désignons par $y_i$ et $z_i$ ($i=1,2$) les points fermés de $Y_i$ ($i=1,2$) correspondant aux idéaux maximaux $\underline{m}$ resp. $\underline{m}'$ de $E$ de hauteur $2$ resp. $1$. Rappelons que les corps résiduels de $\underline{m}$ et $\underline{m}'$ sont isomorphes. On peut donc construire par recollement de $Y_1$ et $Y_2$ un schéma $X$ en identifiant $y_1$ avec $z_2$ et $y_2$ avec $z_1$. On peut alors identifier $Y_1$ et $Y_2$ avec les deux composantes irréductibles de $X$, et alors dans $Z$, $Y_1\cap Y_2=\{y_1,y_2\}$ est de dimension zéro. On voit alors que $X$ est un schéma noethérien, caténaire et équidimensionnel (de dimension $2$). Si (ii bis) était vérifié sur $X$, on aurait
\[
D(y_1)=D(z_1)-1=D(y_2)-1
\]
et d'autre part
\[
D(y_1)=D(z_2)=D(y_2)+1,
\]
donc on trouve une contradiction.

